\chapter{Systemkonzept}
\label{ch:systemkonzept}

\section{Anforderungen}
% TODO: Funktionale Anforderungen (Erkennen, Greifen, Prüfen, Sortieren) und
% nicht-funktionale Anforderungen (Genauigkeit, Taktzeit, Robustheit) als Tabelle.

\section{Hardwareaufbau}
% TODO: Foto/Skizze der Zelle, Roboter, Greifer, Kameras, Beleuchtung, Ablagen.

\section{Softwarearchitektur}
Die Software ist modular aufgebaut. Ein zentraler Orchestrator steuert den Ablauf über
einen Zustandsautomaten und tauscht mit den Modulen ausschließlich definierte Datenobjekte
aus (\texttt{ObjectPose}, \texttt{InspectionResult}, \texttt{RobotTarget}). Jedes Modul
besitzt eine Mock-Implementierung, wodurch die Module unabhängig voneinander entwickelt
und getestet werden können.
% TODO: Architekturdiagramm, Zustandsautomat (Abbildung), Schnittstellentabelle.

\section{Koordinatensysteme}
% TODO: Bild-KS -> Arbeitsbereich-KS (Homographie) -> Roboter-Basis-KS (Starrkörpertransformation).
